% GENERATED FILE. Edit canonical lessons, then run npm run generate:books.
% canonical-source-hash: fnv1a64:d1e84b0d1dcbe0f3
% canonical-lessons: KA-C329-madhyama, KA-C329-vahini, KA-C329-varte, KA-C329-ghosisu, KA-C329-helike

\chapter{The News on Every Channel}
\label{ch:ka-the-news-on-every-channel}
\hlchaptermodality{329}

\begin{chapteropening}
\textbf{By the end of this chapter:} \emph{I can say the media, a TV channel, a news bulletin, to announce and a statement.}

Complete the last lesson of chapter 329: I can say the media, a TV channel, a news bulletin, to announce and a statement.
\end{chapteropening}

\section[\texorpdfstring{\kn{ಮಾಧ್ಯಮ} (mādhyama)}{ಮಾಧ್ಯಮ (mādhyama)}]{\kn{ಮಾಧ್ಯಮ} (mādhyama) — the media; a medium}

\label{lesson:KA-C329-madhyama}



\begin{quote}
Before the new one: say the Kannada for to abscond, then the Kannada for a track.
\end{quote}

\subsection*{You'll want to know: \kn{ಮಾಧ್ಯಮ}}
\textbf{\kn{ಮಾಧ್ಯಮ}} — \emph{mādhyama} — \textquotedblleft{}the media; a medium\textquotedblright{}.

\begin{quote}
\kn{ಮಾಧ್ಯಮಗಳು} \kn{ಆ} \kn{ಘಟನೆಯ} \kn{ಬಗ್ಗೆ} \kn{ದಿನವಿಡೀ} \kn{ಮಾತಾಡಿದವು}.
\end{quote}

\begin{quote}
\emph{mādhyamagaḷu ā ghaṭaneya bagge dinaviḍī mātāḍidavu.}
\end{quote}

\begin{quote}
The media talked about that incident all day.
\end{quote}

Also a medium of teaching: \kn{ಕನ್ನಡ} \kn{ಮಾಧ್ಯಮ} \kn{ಶಾಲೆ}, a Kannada-medium school.

\begin{cousinweb}
\textbf{\kn{ಮಾಧ್ಯಮ}} (\emph{mādhyama}) is borrowed from Sanskrit \emph{mādhyama}, \textquotedblleft{}relating to the middle, central\textquotedblright{} (Monier-Williams), from \emph{madhya}, \textquotedblleft{}middle\textquotedblright{} (CDIAL 9804); the sense \textquotedblleft{}a medium, the media\textquotedblright{} is modern.
\end{cousinweb}

\begin{grammarlens}[title={What you've built}]
One more word for telling what happened.
\end{grammarlens}

\subsection*{Your turn}
\begin{itemize}
\raggedright
  \item \emph{Say it:} \emph{mādhyama}
  \item \emph{Say it:} \emph{mādhyama}, once more
  \item \emph{Say it:} \emph{mādhyamagaḷu ā ghaṭaneya bagge dinaviḍī mātāḍidavu.} — The media talked about that incident all day.
  \item \emph{Recall:} say the Kannada for to abscond, then the Kannada for a track, then say \emph{mādhyama} again
\end{itemize}

\begin{tcolorbox}[breakable,colback=teal!4,colframe=teal!35!black,title={Before you move on}]
What does \kn{ಮಾಧ್ಯಮ} mean? (\textquotedblleft{}The media; a medium\textquotedblright{}.) Say it once more.
\end{tcolorbox}
\section[\texorpdfstring{\kn{ವಾಹಿನಿ} (vāhini)}{ವಾಹಿನಿ (vāhini)}]{\kn{ವಾಹಿನಿ} (vāhini) — a TV channel}

\label{lesson:KA-C329-vahini}



\begin{quote}
Before the new one: say the Kannada for a track, then the Kannada for the media.
\end{quote}

\subsection*{You'll want to know: \kn{ವಾಹಿನಿ}}
\textbf{\kn{ವಾಹಿನಿ}} — \emph{vāhini} — \textquotedblleft{}a TV channel\textquotedblright{}.

\begin{quote}
\kn{ಈ} \kn{ವಾಹಿನಿಯಲ್ಲಿ} \kn{ದಿನವಿಡೀ} \kn{ಸುದ್ದಿ} \kn{ಬರುತ್ತದೆ}.
\end{quote}

\begin{quote}
\emph{ī vāhiniyalli dinaviḍī suddi baruttade.}
\end{quote}

\begin{quote}
This channel has news all day.
\end{quote}

Plural: \kn{ವಾಹಿನಿಗಳು}, \emph{vāhinigaḷu}.

\kn{ಸುದ್ದಿ} \kn{ವಾಹಿನಿ}, a news channel.

\begin{cousinweb}
\textbf{\kn{ವಾಹಿನಿ}} (\emph{vāhini}) is borrowed from Sanskrit \emph{vāhinī}, \textquotedblleft{}an army\textquotedblright{}, and in later use \textquotedblleft{}a river\textquotedblright{}, that which flows and carries (Monier-Williams); the sense \textquotedblleft{}a channel\textquotedblright{} is modern.
\end{cousinweb}

\begin{grammarlens}[title={What you've built}]
One more word for telling what happened.
\end{grammarlens}

\subsection*{Your turn}
\begin{itemize}
\raggedright
  \item \emph{Say it:} \emph{vāhini}
  \item \emph{Say it:} \emph{vāhini}, once more
  \item \emph{Say it:} \emph{ī vāhiniyalli dinaviḍī suddi baruttade.} — This channel has news all day.
  \item \emph{Recall:} say the Kannada for a track, then the Kannada for the media, then say \emph{vāhini} again
\end{itemize}

\begin{tcolorbox}[breakable,colback=teal!4,colframe=teal!35!black,title={Before you move on}]
What does \kn{ವಾಹಿನಿ} mean? (\textquotedblleft{}A TV channel\textquotedblright{}.) Say it once more.
\end{tcolorbox}
\section[\texorpdfstring{\kn{ವಾರ್ತೆ} (vārte)}{ವಾರ್ತೆ (vārte)}]{\kn{ವಾರ್ತೆ} (vārte) — a news bulletin; news}

\label{lesson:KA-C329-varte}



\begin{quote}
Before the new one: say the Kannada for the media, then the Kannada for a TV channel.
\end{quote}

\subsection*{You'll want to know: \kn{ವಾರ್ತೆ}}
\textbf{\kn{ವಾರ್ತೆ}} — \emph{vārte} — \textquotedblleft{}a news bulletin; news\textquotedblright{}.

\begin{quote}
\kn{ಅಜ್ಜ} \kn{ಪ್ರತಿದಿನ} \kn{ರಾತ್ರಿ} \kn{ಎಂಟು} \kn{ಗಂಟೆಯ} \kn{ವಾರ್ತೆ} \kn{ಕೇಳುತ್ತಾರೆ}.
\end{quote}

\begin{quote}
\emph{ajja pratidina rātri eṇṭu gaṇṭeya vārte kēḷuttāre.}
\end{quote}

\begin{quote}
Grandfather listens to the eight o'clock news bulletin every night.
\end{quote}

\kn{ವಾರ್ತೆ} is the formal word on radio and television; in talk people say \kn{ಸುದ್ದಿ}.

\begin{cousinweb}
\textbf{\kn{ವಾರ್ತೆ}} (\emph{vārte}) is borrowed from Sanskrit \emph{vārttā}, \textquotedblleft{}an account of what has happened, tidings, news\textquotedblright{} (Monier-Williams).
\end{cousinweb}

\begin{grammarlens}[title={What you've built}]
One more word for telling what happened.
\end{grammarlens}

\subsection*{Your turn}
\begin{itemize}
\raggedright
  \item \emph{Say it:} \emph{vārte}
  \item \emph{Say it:} \emph{vārte}, once more
  \item \emph{Say it:} \emph{ajja pratidina rātri eṇṭu gaṇṭeya vārte kēḷuttāre.} — Grandfather listens to the eight o'clock news bulletin every night.
  \item \emph{Recall:} say the Kannada for the media, then the Kannada for a TV channel, then say \emph{vārte} again
\end{itemize}

\begin{tcolorbox}[breakable,colback=teal!4,colframe=teal!35!black,title={Before you move on}]
What does \kn{ವಾರ್ತೆ} mean? (\textquotedblleft{}A news bulletin; news\textquotedblright{}.) Say it once more.
\end{tcolorbox}
\section[\texorpdfstring{\kn{ಘೋಷಿಸು} (ghōṣisu)}{ಘೋಷಿಸು (ghōṣisu)}]{\kn{ಘೋಷಿಸು} (ghōṣisu) — to announce; to declare}

\label{lesson:KA-C329-ghosisu}



\begin{quote}
Before the new one: say the Kannada for a TV channel, then the Kannada for a news bulletin.
\end{quote}

\subsection*{You'll want to know: \kn{ಘೋಷಿಸು}}
\textbf{\kn{ಘೋಷಿಸು}} — \emph{ghōṣisu} — \textquotedblleft{}to announce; to declare\textquotedblright{}.

\begin{quote}
\kn{ಸರ್ಕಾರ} \kn{ನಾಳೆ} \kn{ರಜೆ} \kn{ಅಂತ} \kn{ಘೋಷಿಸಿತು}.
\end{quote}

\begin{quote}
\emph{sarkāra nāḷe raje anta ghōṣisitu.}
\end{quote}

\begin{quote}
The government announced that tomorrow is a holiday.
\end{quote}

Past (I …): \kn{ಘೋಷಿಸಿದೆ}, \emph{ghōṣiside}.

What is announced comes before \kn{ಎಂದು} or \kn{ಅಂತ}: \kn{ರಜೆ} \kn{ಅಂತ} \kn{ಘೋಷಿಸು}.

\begin{cousinweb}
\textbf{\kn{ಘೋಷಿಸು}} (\emph{ghōṣisu}) is built on Sanskrit \emph{ghoṣa}, \textquotedblleft{}a noise, the cry of a crowd\textquotedblright{}, whose derivative \emph{ghoṣaṇā} means \textquotedblleft{}a proclamation\textquotedblright{} (Monier-Williams), with the Kannada verb ending -\emph{isu}.
\end{cousinweb}

\begin{grammarlens}[title={What you've built}]
One more word for telling what happened.
\end{grammarlens}

\subsection*{Your turn}
\begin{itemize}
\raggedright
  \item \emph{Say it:} \emph{ghōṣisu}
  \item \emph{Say it:} \emph{ghōṣisu}, once more
  \item \emph{Say it:} \emph{sarkāra nāḷe raje anta ghōṣisitu.} — The government announced that tomorrow is a holiday.
  \item \emph{Recall:} say the Kannada for a TV channel, then the Kannada for a news bulletin, then say \emph{ghōṣisu} again
\end{itemize}

\begin{tcolorbox}[breakable,colback=teal!4,colframe=teal!35!black,title={Before you move on}]
What does \kn{ಘೋಷಿಸು} mean? (\textquotedblleft{}To announce; to declare\textquotedblright{}.) Say it once more.
\end{tcolorbox}
\section[\texorpdfstring{\kn{ಹೇಳಿಕೆ} (hēḷike)}{ಹೇಳಿಕೆ (hēḷike)}]{\kn{ಹೇಳಿಕೆ} (hēḷike) — a statement}

\label{lesson:KA-C329-helike}



\begin{quote}
Before the new one: say the Kannada for a news bulletin, then the Kannada for to announce.
\end{quote}

\subsection*{You'll want to know: \kn{ಹೇಳಿಕೆ}}
\textbf{\kn{ಹೇಳಿಕೆ}} — \emph{hēḷike} — \textquotedblleft{}a statement\textquotedblright{}.

\begin{quote}
\kn{ಪೊಲೀಸರ} \kn{ಹೇಳಿಕೆ} \kn{ಎಲ್ಲ} \kn{ಪತ್ರಿಕೆಗಳಲ್ಲಿ} \kn{ಬಂತು}.
\end{quote}

\begin{quote}
\emph{polīsara hēḷike ella patrikegaḷalli bantu.}
\end{quote}

\begin{quote}
The police statement appeared in every newspaper.
\end{quote}

\kn{ಹೇಳು}, to say, with -\kn{ಇಕೆ}, which makes a noun of a verb: what is said.

\begin{cousinweb}
\textbf{\kn{ಹೇಳಿಕೆ}} (\emph{hēḷike}) is a native Kannada noun made from the verb \emph{hēḷu}, \textquotedblleft{}to say\textquotedblright{} (DEDR 4430), with the ending -\emph{ike} that turns verbs into nouns.
\end{cousinweb}

\begin{grammarlens}[title={What you've built}]
One more word for telling what happened.
\end{grammarlens}

\subsection*{Your turn}
\begin{itemize}
\raggedright
  \item \emph{Say it:} \emph{hēḷike}
  \item \emph{Say it:} \emph{hēḷike}, once more
  \item \emph{Say it:} \emph{polīsara hēḷike ella patrikegaḷalli bantu.} — The police statement appeared in every newspaper.
  \item \emph{Recall:} say the Kannada for a news bulletin, then the Kannada for to announce, then say \emph{hēḷike} again
\end{itemize}

\begin{tcolorbox}[breakable,colback=teal!4,colframe=teal!35!black,title={Before you move on}]
What does \kn{ಹೇಳಿಕೆ} mean? (\textquotedblleft{}A statement\textquotedblright{}.) Say it once more.
\end{tcolorbox}
